\section{第一阶段：利用现有基础模型改造机器人三板块}

上一章提出两阶段路线，本章看第一阶段：把现有 LLM、VLM 和 Code LM 放进机器人系统。传统机器人三板块是 Planning、Perception、Actuation。大模型首先替代规划，因为语言模型擅长把任务拆成步骤；随后 VLM 替代感知，因为它能把视觉场景和语言任务连接起来；再进一步，Code LM 试图把执行也自动化。

\subsection{SayCan：我能做，而不是我说做}

SayCan 的核心问题是：语言模型会提出很多动作，但机器人未必能在当前环境中做到。它将语言模型的任务分解和机器人 affordance 结合起来，既考虑“语言上应该做什么”，也考虑“机器人实际上能不能做”。这是把 LLM 接到机器人规划上的开创性工作。

\begin{figure}[H]
\centering
\includegraphics[width=0.96\textwidth]{ch07-01537.jpg}
\caption{SayCan：Query LLM to rank action primitives，再用 affordance function 判断可行性。投屏帧，约 00:25:37。}
\end{figure}

\begin{knowledgebox}{读图：LLM 给目标，affordance 给约束}
图中左侧是语言模型对动作的排序，右侧是机器人在现实环境中的可供性判断。读这张图时要注意，SayCan 没有让语言模型直接控制机器人，而是在语言计划和物理可行性之间加了一道约束。
\end{knowledgebox}

\begin{warningbox}{语言计划不能直接等于机器人动作}
机器人需要考虑抓取能力、空间位置、物体状态和执行失败。LLM 如果只输出“拿起杯子”，并不知道杯子是否可达、夹爪是否合适、环境是否允许。SayCan 的价值就是把语言意图落地到 affordance。
\end{warningbox}

\subsection{Inner Monologue、DoReMi 与 VoxPoser}

本节看三篇继续把基础模型用于机器人推理的工作。Inner Monologue 强调机器人执行过程中需要语言化的反馈和内心独白，把环境反馈重新纳入规划。DoReMi 关注计划和执行不一致时如何检测并恢复。VoxPoser 则用语言模型构造 3D value maps，把语言目标转成可组合的空间价值场，用于 manipulation。

\begin{figure}[H]
\centering
\includegraphics[width=0.90\textwidth]{ch08-01649.jpg}
\caption{Inner Monologue：通过语言模型规划和环境反馈形成具身推理循环。投屏帧，约 00:27:29。}
\end{figure}

\begin{knowledgebox}{读图：机器人需要执行后的自我叙述}
Inner Monologue 的图把 Action、Environmental Feedback、Corrective Action 连成闭环。它不是简单让机器人“会说话”，而是让执行结果回到语言规划过程，形成可修正的中间状态。
\end{knowledgebox}

\begin{figure}[H]
\centering
\includegraphics[width=0.94\textwidth]{ch09-01737.jpg}
\caption{DoReMi：检测并恢复 plan-execution misalignment。投屏帧，约 00:28:57。}
\end{figure}

\begin{knowledgebox}{读图：错配检测是具身智能的现实问题}
DoReMi 处理的是规划与执行错配：计划里说 A，执行中环境变成 B，或者机器人没有按预期完成。它提醒我们，机器人系统不是一次生成计划就结束，而是要持续检测、纠错和恢复。
\end{knowledgebox}

\begin{figure}[H]
\centering
\includegraphics[width=0.94\textwidth]{ch11-01956.jpg}
\caption{VoxPoser：用语言模型组合 3D value maps 以指导机器人操作。投屏帧，约 00:32:36。}
\end{figure}

\begin{knowledgebox}{读图：从语言目标到 3D 操作场}
VoxPoser 图中把语言指令转成空间中的 value map。读图时要看“可组合”二字：不同约束可以叠加成一个操作目标，例如接近物体、避开障碍、保持姿态等。
\end{knowledgebox}

\subsection{从 Code LM 到真正的机器人模型}

Code LM 替代 Actuation 的想法，是让模型写机器人代码或调用底层 API。但陈建宇指出，这仍然不是完整的机器人 foundation model。它更多是在现有工具和现有控制器上做编排；真正的 VLA 需要模型本身具备从视觉和语言直接输出动作的能力，而不是只会写调用代码。

\begin{importantbox}{第一阶段的边界}
Leveraging foundation models 能快速增强规划、感知和工具调用，但它本质上仍是模块化系统。通用机器人最终需要在机器人数据上训练能够处理 action 的模型。
\end{importantbox}

\subsection{本章小结}

第一阶段的贡献，是证明 LLM/VLM 能帮助机器人理解任务、计划行动、检测错配和构造空间目标。但这些方法仍依赖传统机器人模块，尚未真正形成端到端 VLA。

